\documentclass[11pt]{article}
\usepackage[margin=1in]{geometry}
\usepackage{amsmath,amssymb,amsthm,mathtools}
\newtheorem{proposition}{Proposition}
\newtheorem{lemma}{Lemma}
\newtheorem{conjecture}{Conjecture}
\theoremstyle{definition}
\newtheorem{target}{Target}
\newtheorem{remark}{Remark}
\newcommand{\R}{\mathbb{R}}
\newcommand{\E}{\mathbb{E}}
\newcommand{\Prob}{\mathbb{P}}

\title{Theory targets after the S3 wide-band sweep\\
\large candidate statements for the author to prove, correct, or extend}
\author{Working notes; nothing here goes into the paper until proven and verified}
\date{2 July 2026}
\begin{document}
\maketitle

\paragraph{Context.} The S3 sweep (AT strength $\varepsilon\in\{2,4,8,16\}/255$ $\times$ dose
$\{0,1\mathrm M\}$ $\times$ 2 seeds, $n=16$ cells, threat-matched AutoAttack at the training
$\varepsilon$) produced two facts that currently have no theory layer in the paper:
\begin{itemize}
\item[(F1)] the \emph{budget-normalized} ratio $\eta/(L_1\varepsilon)$ predicts AutoAttack robust
accuracy across the whole sweep (Pearson $+0.94$, Spearman $+0.94$), while the raw ratio
$\eta/L_1$ anti-correlates across budgets ($-0.95$) because a harder budget lowers AA while raising
the raw margin;
\item[(F2)] at matched $\eta/L_1$ the $+1$M-synthetic models sit $+0.10$ AA above the real-only
models ($t\approx+5$), the offset tracks a smaller train--test robust gap ($-0.45$), and it vanishes
at $\varepsilon=16/255$ where there is no robust overfitting left to prevent.
\end{itemize}
Targets are ordered by (value $\times$ tractability). T1--T2 are new and, I believe, within reach;
T3 answers a specific referee question; T4--T5 are the standing harder items.

\section*{T1. The per-sample certificate count is the budget-normalized law}

\paragraph{Claim to prove (exact first-order layer).}
For the $\ell_\infty$ threat at budget $\varepsilon$, define the per-sample first-order radius
$r_1(x)=M(x)/\|\nabla M(x)\|_1$ (the $\ell_\infty$ Hein radius; dual norm $\ell_1$).

\begin{proposition}[Candidate: survival $=$ certificate, affine case]\label{prop:count}
If $M$ is affine on the ball $B_\infty(x,\varepsilon)$ and $M(x)>0$, then
$x$ survives every perturbation of $\ell_\infty$ norm $\le\varepsilon$ if and only if
$r_1(x)\ge\varepsilon$. Consequently, for a model that is affine on every such ball,
\[
\mathrm{RobAcc}(\varepsilon)\;=\;\Prob_x\big(r_1(x)\ge\varepsilon,\ M(x)>0\big)
\quad\text{(the \emph{certified fraction}).}
\]
\end{proposition}
\noindent\emph{Proof route.} One direction is the Hein/H\"older certificate; the converse takes
$\delta=-\varepsilon\,\mathrm{sign}(\nabla M(x))$, which attains the dual norm exactly in
$\ell_\infty$. Should be short and clean. The paper currently has no statement at the
\emph{robust-accuracy} (population count) level; this is the missing bridge between the per-sample
certificate (\S5 ``per sample'') and the cross-model law.

\paragraph{Claim to prove (why the \emph{mean} ratio works).}
The empirical law (F1) uses the ratio of means $\eta/L_1=\E[M]/\E\|\nabla M\|_1$, not the certified
fraction. Candidate explanation: a scale-family hypothesis.

\begin{conjecture}[Scale family]\label{conj:scale}
Across the trained models of one family (fixed architecture and data, varying AT strength and dose),
the per-sample radius distribution is approximately a scale family:
$r_1(x)\overset d=\mu\,Z$ with $\mu=\mu(\text{model})$ a scale parameter and $Z$ a
model-independent shape variable. Then
$\mathrm{RobAcc}(\varepsilon)=\Prob(Z\ge\varepsilon/\mu)=F_Z(\mu/\varepsilon)$ is a fixed monotone
function of $\mu/\varepsilon$, and any scale proxy (the mean ratio, the median $r_1$) collapses the
sweep onto one curve. This would \emph{derive} the budget-normalized law and predict that
$\mathrm{AA}$ vs $\eta/(L_1\varepsilon)$ is not just correlated but a single monotone curve.
\end{conjecture}
\noindent\emph{What to do.} (i) I can test the shape-invariance empirically from the saved S3
checkpoints (per-sample $r_1$ histograms, normalized by their median, overlaid across cells); (ii)
if it holds, the paper statement is ``the sweep collapses onto $F_Z$''; a distribution-free
one-sided version (Markov/Chebyshev on $1/r_1$) gives bounds without the hypothesis. Author call on
how much to formalize.

\section*{T2. Is the data-axis offset a tail effect? (highest value)}

\paragraph{The point.} By Proposition~\ref{prop:count}, robust accuracy is a \emph{tail
functional} of the $r_1$ distribution, while $\eta/L_1$ is a \emph{ratio of means}. Two models can
match in mean ratio and differ in the lower tail. Candidate: the $+0.10$ offset (F2) is exactly
this — the synthetic-data model has a thinner lower tail (fewer barely-non-robust test points) at
matched mean, because robust overfitting manufactures a heavy lower tail on \emph{test} points
(train points get memorized margins, test points near memorized islands lose theirs).

\begin{conjecture}[Tail, not mean]\label{conj:tail}
Conditioning on the certified fraction $\Prob_x(r_1(x)\ge\varepsilon)$ instead of the mean ratio
absorbs the dose coefficient: in the OLS of AA on $\{$certified fraction, dose$\}$ the dose beta is
statistically zero. Equivalently, the data intervention improves robustness \emph{through} the
$r_1$ tail, and the ``generalization offset'' is the mean being an insufficient statistic, not a
beyond-certificate effect.
\end{conjecture}
\noindent\emph{What to do.} (i) [cheap, I run it] recompute per-sample $r_1$ on the report split
from the 16 saved S3 checkpoints; regress AA on certified-fraction $+$ dose; also check the
variance/quantiles of $r_1$ at matched mean. (ii) [theory, author] a small lemma making the
mechanism precise, e.g.\ in the vicinal two-mode model: coverage reduces the \emph{across-test-point
variance} of $r_1$ at a given mean (the rough mode contributes idiosyncratic per-point sensitivity;
suppressing it homogenizes $r_1$), hence raises any lower-tail functional at matched mean. If
(i) confirms and (ii) is proven, the offset Limitation becomes a theorem-backed refinement and the
vicinal appendix acquires its missing dynamical consequence.

\section*{T3. A certifiable co-Lipschitz $\alpha$ beyond the linear case}

Referee question (theory reviewer, Q1): exhibit a nonlinear, non-toy class where the bracket's
upper arm is informative, i.e.\ $\alpha$ certified \emph{independently} of $r_2$.
Route (extends E4 of the earlier audit): scores $g=h\circ\Psi$ with $h$ monotone and $\Psi$ a
quadratic invariant. On a cone $\mathcal C$ bounded away from the phase-blind directions, the
power-spectrum map satisfies an explicit two-sided bound
$c_-(\mathcal C)\|x-y\|\le\|\Psi(x)-\Psi(y)\|\le 2B\|x-y\|$ along radial-ish paths; composing with
$h'\ge h_->0$ gives $\alpha=h_-c_-$ certified from the feature structure alone. The deliverable is
a proposition with the explicit cone constant $c_-$ and a numerical check on the surrogate we
already verified ($\kappa\approx1.6$--$2.4$). Medium difficulty; answers the referee directly.

\section*{T4. A dynamical vicinal-coverage statement (toy)}

The appendix's static path result explains coverage; the verified mechanism is dynamical
(robust-overfitting delay). Candidate toy bridge, in the two-mode model of thm:vic-twomode: give the
rough coordinate an epoch-driven memorization drift (fitting label noise on the empirical measure,
$\dot a_r\propto+\sigma^2$ on train, harming test), and show the coverage term $tq$ multiplies its
decay so the test-ratio collapse time scales like $t_{\mathrm{collapse}}\sim(1+tq)/\sigma^2$.
Even at this cartoon level it would let the appendix say ``coverage delays the collapse'' as a
statement rather than a reading. Medium-hard; optional but would upgrade App.~vicinal.

\section*{T5. Rich-regime dynamics selection (standing open problem)}

Unchanged: whether GD on generic orbit data selects an invariant / lower-realized-$L$ optimum among
the equal-variation-norm max-margin interpolants (the non-strict-convexity gap in
cor:no-gap-var). This is the wall of Frei/Min/Li and stays the paper's stated open problem; no new
angle from S3.

\begin{remark}[Division of labor]
T2(i) and the T1 shape-check are empirical and cheap: I run them from the saved checkpoints without
touching the paper. T1's proposition, T2's lemma, T3's proposition, and T4 are author theory.
Nothing enters the paper until proven by the author and verified firsthand, per the standing rule.
\end{remark}

\end{document}
